\documentclass[aspectratio=169,11pt]{beamer}
\usepackage[utf8]{inputenc}
\usepackage[T1]{fontenc}
\usepackage[spanish]{babel}
\usepackage{lmodern}
\usepackage{booktabs}
\usepackage{tabularx}
\usepackage{hyperref}
\usepackage{enumitem}

\usetheme{Madrid}
\usecolortheme{default}
\setbeamertemplate{navigation symbols}{}
\setbeamertemplate{footline}[frame number]

\title{IA agéntica para la creación de contenido académico}
\subtitle{Sesión 7 --- Artículos I: estructura, borrador asistido y citas}
\author{Juliho Castillo}
\institute{CADi 40\,h · Escuela de Ingeniería y Ciencias\\
Tecnológico de Monterrey}
\date{Clave sugerida: \texttt{CADI-EIC-IA-AGENTICA-01}}

\begin{document}

\begin{frame}
\titlepage
\end{frame}

\begin{frame}{Agenda (3\,h) y producto E7}
\begin{tabularx}{\textwidth}{@{}lclX@{}}
\toprule
Bloque & Tiempo & Contenido \\
\midrule
Marco artículo & 15 min & IMRyD, audiencia, contribución \\
Demo borrador + citas & 25 min & No fabricar referencias \\
Taller estructura/borrador & 85 min & Micro-sección o guía \\
Lista de verificación & 25 min & Citas y datos \\
Cierre & 20 min & Iniciar agent log; TI \\
\bottomrule
\end{tabularx}
\vspace{0.5em}
\textbf{E7:} estructura + borrador con citas/datos a verificar; inicio de agent log.
\end{frame}

\begin{frame}{Resultados de aprendizaje (sesión)}
Al finalizar, usted podrá:
\begin{enumerate}
\item Estructurar un \textbf{micro-artículo/sección} (p.\ ej.\ IMRyD) o guía breve
\item Producir un \textbf{borrador asistido} con marcas \texttt{[VERIFICAR]}
\item Elaborar una \textbf{lista de citas/datos} a verificar antes de uso público
\item Iniciar el \textbf{agent log} del pipeline completo
\end{enumerate}
\end{frame}

\begin{frame}{Marco: el microartefacto puede ser\ldots}
\begin{itemize}
\item Sección tipo IMRyD (o variante disciplinar)
\item Guía de laboratorio breve
\item Consolidación de libro/deck en prosa académica corta
\end{itemize}
\begin{block}{Alcance CADi}
Priorice calidad de $\sim$800--1200 palabras (o guía equivalente), no un paper completo.
\end{block}
\end{frame}

\begin{frame}{Estructuras académicas (IMRyD y variantes)}
\begin{tabularx}{\textwidth}{@{}lX@{}}
\toprule
Bloque & Pregunta que responde \\
\midrule
Introducción & ¿Qué problema y por qué importa? \\
Método / enfoque & ¿Cómo se aborda? \\
Resultados / desarrollo & ¿Qué se obtuvo o explica? \\
Discusión / cierre & ¿Qué implica? ¿Límites? \\
\bottomrule
\end{tabularx}
\vspace{0.5em}
Adapte a guías: objetivo, materiales, procedimiento, riesgos, cierre.
\end{frame}

\begin{frame}{Audiencia y contribución}
Antes de redactar:
\begin{itemize}
\item ¿Quién lee? (estudiantes, pares, taller)
\item ¿Qué aporta que no sea genérico?
\item ¿Qué afirmaciones requieren evidencia?
\item ¿Qué queda fuera del microartefacto?
\end{itemize}
\end{frame}

\begin{frame}{Demo --- borrador + \texttt{[VERIFICAR]}}
En vivo:
\begin{enumerate}
\item Generar párrafo con restricción ``no inventar citas''
\item Detectar cita fantasma o DOI dudoso
\item Marcar \texttt{[VERIFICAR]} o eliminar
\item Humano corrige con fuente real o quita la afirmación
\end{enumerate}
\begin{alertblock}{Regla de oro}
La IA propone; la persona verifica. \textbf{Nunca} se fabrican citas.
\end{alertblock}
\end{frame}

\begin{frame}{Citas: verificación humana obligatoria}
\begin{itemize}
\item Cada referencia debe existir y decir lo que se afirma
\item Sin DOI/URL verificable $\rightarrow$ no citar o marcar pendiente
\item Datos/tablas: fuente o \texttt{[VERIFICAR]}
\item ``Se ve académico'' no es criterio de aceptación
\end{itemize}
\begin{block}{CEDDIE}
Política de integridad / uso de IA: verificar por edición; no improvisar normativa.
\end{block}
\end{frame}

\begin{frame}{Plantilla de lista de citas/datos}
Para cada ítem:
\begin{enumerate}
\item Afirmación en el texto
\item Cita/dato propuesto
\item Estado: verificado / pendiente / eliminado
\item Método: DOI, base de datos, libro, experta/o
\item Responsable humano
\end{enumerate}
\end{frame}

\begin{frame}{Inicio del agent log}
Documente desde hoy (se cierra en sesión 8):
\begin{itemize}
\item Herramienta / modelo
\item Rol del agente y brief/prompt
\item Qué generó la IA
\item Qué cambió o rechazó la persona
\item Pendientes de verificación
\end{itemize}
\texttt{[VERIFICAR]} es hábito profesional, no debilidad.
\end{frame}

\begin{frame}{Roles en el pipeline de artículo}
\begin{itemize}
\item Planificador: estructura IMRyD / guía
\item Redactor: borrador por bloques cortos
\item Verificador de citas: lista de pendientes
\item Humano: decide qué se publica
\end{itemize}
\end{frame}

\begin{frame}{Actividad de taller ($\sim$85 min + lista)}
\textbf{Entregable E7:}
\begin{enumerate}
\item Estructura del micro-texto o guía
\item Borrador asistido con \texttt{[VERIFICAR]} visibles
\item Lista de citas/datos a verificar + método
\item Agent log iniciado
\end{enumerate}
Nombre: \texttt{Apellido\_Sesion7\_borrador-citas} (+ log)
\end{frame}

\begin{frame}{Rúbrica rápida (E7)}
\begin{tabularx}{\textwidth}{@{}lXX@{}}
\toprule
Criterio & Bien & A mejorar \\
\midrule
Estructura & Bloques claros & Texto amorfo \\
Integridad & Sin citas inventadas & Fantasmas \\
HITL & Lista + \texttt{[VERIFICAR]} & Aceptación pasiva \\
Log & Iniciado y usable & Ausente \\
\bottomrule
\end{tabularx}
\end{frame}

\begin{frame}{Anti-patrones}
\begin{itemize}
\item Bibliografía inventada ``para que se vea serio''
\item Pegar el borrador al LMS sin lista de verificación
\item Agent log vacío (``ya lo haré al final'')
\item Afirmaciones empíricas sin fuente ni marca
\end{itemize}
\end{frame}

\begin{frame}{Takeaways de la sesión}
\begin{enumerate}
\item Estructura primero; prosa después
\item Citas: verificación humana obligatoria
\item \texttt{[VERIFICAR]} es profesionalismo
\item El agent log empieza hoy
\end{enumerate}
\end{frame}

\begin{frame}{Trabajo independiente --- bloques C + D}
Antes de la sesión 8:
\begin{itemize}
\item Completar borrador (C)
\item Avanzar agent log (D)
\item Lectura corta del bloque A si aplica
\item Traer microartefacto casi cerrado para revisión
\end{itemize}
Parte de las \textbf{16\,h} independientes.
\end{frame}

\begin{frame}{Prep sesión 8 --- Artículos II y cierre}
Última sesión en aula (3\,h):
\begin{itemize}
\item Agentes de revisión (claridad, rigor, sesgos)
\item Cierre del microartefacto + checklist
\item Showcase de aprendizajes
\item Cierre del CADi
\end{itemize}
\end{frame}

\begin{frame}{Gracias}
\begin{center}
{\Large Juliho Castillo}\\[0.5em]
Escuela de Ingeniería y Ciencias\\[1em]
CADi · IA agéntica para la creación de contenido académico\\[1.5em]
{\small \emph{Verificar políticas de integridad y uso de IA con CEDDIE antes de cada edición.}}
\end{center}
\end{frame}

\end{document}
